\section{The Guardian of the Threshold}

One of the biggest problems with modern political entities is the question of cybernetics, or ``the helmsmen'' : who watches the Watchers? Who guards the guardians? This problem is not very effectively answered today; if there was a famine in the year 1200, you knew\footnote{\url{https://en.wikipedia.org/wiki/Jacquerie}} where the nobleman who was feasting on venison lived, what his name was, \& whether or not he had contributed to the problem. Likewise, the king oversaw them in their turn, \& perhaps an emperor looked over the king.

Modern man doesn't have an answer to ``who the guardians'' are, in either the terrestial life, or the life-to-come. Only the ``Free Market'' regulates such matters.

Esotericism (however) has a very precise, and relatively homogeneous answer to this problem, which goes even beyond political solutions. The ``Guardian of the Threshold\footnote{\url{https://en.wikipedia.org/wiki/Guardian\_of\_the\_Threshold}}'' is an entity personal to each human being, and consists of the sum total ``bad'' or negative acts of that being -- this ``doppelganger'' constrains, judges, and (perhaps) terrifies that being, in this life \& the next.

Tomasso Palamidessi articulated the strongest doctrines concerning this entity, but this idea is by no means peculiar to merely him. Papus mentions\footnote{\url{http://www.sacred-texts.com/tarot/tob/index.htm}} the movement through the Zodiac houses which every being undertakes as part of its journey, and one of these entities or houses is ``the guardian''. Likewise, Max Heindel speaks of the way the thought/mental, vital, and spiritual bodies interact in a way that is suggestive of the idea that we are surrounded by our own cosmic creations:

\begin{quotex}
Repulsion is the centrifugal force and if that is aroused by the thought there will be a struggle between the spiritual force ()the will of the man) within the though-form, and the desire body. This is the battle between conscience and desire, the higher and the lower nature. The spiritual desire0stuff needed to manipulate the brain and muscles. The force of repulsion will endeavor to scatter the appropriated material and oust the thought. If the spiritual energy is strong it may force its way through to brain center and hold its clothing of desire-stuff while manipulating the vital force, thus compelling action, and will then leave upon the memory a vivid impression of the struggle and the victory. If the spiritual energy is exhausted before action has resulted, it will be overcome by the force of Repulsion , and will be stored in the memory, as are all other thought forms when they have expended their energy.
\flright{\emph{Rosicrucian Cosmo-Conception}, p. 26}
\end{quotex}


These facts or truths may be behind the old medieval belief that demons were attempting to pluck man off the ladder, or carry his soul to hell when he died. Heindel goes on\footnote{\url{http://www.rosicrucian.com/rcc/rcceng00.htm}} to speak of ``ripe Fate'' which not even the Lords of Destiny will enable man to ``shirk''.

Franz Bardon also makes mention\footnote{\url{http://www.abardoncompanion.com/IIH-Step1.html}} in his magical system of mental ``larva'' which are generated by frustrated or misused desire.

Valentine Tomberg speaks of the ``Guardian'' in some detail in Chapter XVIII of \emph{Meditations on the Tarot,} particularly in relation to conscience.

\begin{quotex}
The guardian of the threshold spoken of in the Hermetic tradition is the great judge, charged with preserving the equilibrium of that which is above and that which is below. The traditional iconography of the Church represents him with a sword and balance. The sword is his vivifying and healing action, giving courage and humility to the soul which hungers and thirsts for the depths, and the balance is his action of presenting the precise account of what must be paid in order to have the right to go further.
\end{quotex}


Thus, we see that behind the old medieval traditions of visions, temptations,\footnote{\url{http://classicalchristianity.com/2011/10/29/the-incredible-vision-of-st-drythelm/}} sin, the devil, etc. lurked a very prescient sense of the true reality of man's situation, which is always perilous, always blessed. The modern world derides ``shame cultures\footnote{\url{http://en.wikipedia.org/wiki/Shame\_society}}'' as inherently backwards and perverse, yet it has nothing better to put in its place, nor has it managed to do away with ``shame'', if for no other reason than that it can pass all the progressive legislation it wishes to, it cannot abolish the reality of the Guardian of the Threshold. This Guardian waits for every man, and for those who cannot don the armor of St. George whereby to triumph over the old serpent, a heavy price will be paid.

Such a spiritual truth is surely behind the old Scripture verse : ``It is appointed unto man once to die, and after that, the judgement''. Even re-incarnation or the Lords of Destiny will not save a man from paying the full and terrible price for those deeds, for ``re-birth'' is not an erasure of karma, but rather a further ripening of it, a passing to an new \& developmental state.

There is, in the Christian religion, one possible way out of karma\footnote{\url{https://en.wikipedia.org/wiki/Body\_of\_resurrection}}, but it is not a release from its burden, but rather a servant-like assumption of it in a state of pleromic Love, the highest state of Being identical to the Lord Himself. This is known as ``forgiveness'' and ``the new birth'', and rather than being a blank check for modern people to silence their inner ``guardians'' of conscience, it represents a higher path which allows a victory where perhaps none is possible. This can occur at the whim of the Lord, and therefore, is not ultimately subject to earthly protocol or guarantees, either in its favor (\emph{extra ecclesiam non est salus}) or (as the case today is) against it as a presumption upon grace in the form of graceless and shameless culture.


\flrightit{Posted on 2011-12-30 by Logres}

\begin{center}* * *\end{center}

\begin{footnotesize}
\begin{sffamily}
\texttt{Logres on 2011-12-30 at 22:45 said:}

I should note that this ``entity'' is MORE than merely the lower doppelganger or judge -- the descriptions given are from various angles, involving various aspects of something that actually exists in the depths and heights of a person's being, vertically, up or down.

\hfill

\texttt{Attila on 2011-12-31 at 22:05 said:}

All very interesting -- but when push comes to shove -- we all need a method of realizing our true nature within, through and in The Whole. One of the most beautiful ways is to be found in the Sufi tradition --- as shown here:

\url{http://www.youtube.com/watch?v=\_SEs\_CDciKk}


\hfill

\texttt{logres on 2012-01-01 at 19:55 said:}

Looks good, but I don't speak Arabic. Can only intuit. The voice is hypnotic, however.

\hfill

\texttt{Charlotte Cowell on 2012-01-02 at 05:58 said:}

This is one of my favourite Arabic songs, and being devoted to a beloved I think we can give it some Sufi credentials :-)

\url{http://www.youtube.com/watch?v=YTZO1DNPqas\&feature=related}


\hfill

\texttt{Caleb Cooper on 2012-01-03 at 02:45 said:}

Brings back memories. When I started waking up I began perceiving demonic presences all over the place. I'm pretty sure now that at least most of them were lower aspects of my soul that I found repugnant and walled off. Fortunately I've learned to be a lot more gentle with them and myself.

At my first Shamanic training 8 months ago (foundation for shamanic studies) on my first journey, I was attacked by my totem at the time, a black panther. Which was very good, as we were no longer suited for each other; I survived my adolescence thanks to its strength, but it also represented the viciousness with which I defended my desire to be solitary.

What I needed though was to start opening up my heart horizontally, so I subdued the totem and released it. By the end of the workshop God provided me with a perfect primary totem for what I needed to learn, a quail (before you would fly, learn to walk grounded with more humility).

Recently, I had a pretty bad experience when I decided to try and pass the Trail of Fire. I was incorporating a Tibetan Technique of coming into resonance with a god and imprinting it upon the corresponding chakra (actually, the Tibetans \_absorb\_ the god into themselves, but I can only manage temporary merging, and that's on a \_good\_ day).

I didn't maintain the purity needed though and ended up somewhere bad, getting attacked by what a demonic aspect of Shiva might look like. It doesn't do any good ``killing'' something like that, because it just dissolves and covers you in it's crap, clogging your chakras for days. Maybe it was an aspect of myself, but it felt like an ``other'', not that I'm trusting that perception.

\hfill

\texttt{Charlotte Cowell on 2012-01-03 at 07:14 said:}

Caleb, I also underwent Shamanic training for many years, it was the first thing I really learned, which is appropriate as it is bound to come naturally to a human creature. Interestingly in the context of your own story, the occasion of my initiation into universal Christianity was the time I first came into direct contact/awareness with the animal totem, which in my case was a leopard. There is something highly symbolic about this animal, which I suppose you have already looked into. Almost immediately afterwards I found myself talking to and with birds, so again there is a parallel. In particular owls, wrens and eventually an Eagle, but this was more of a guide. Unlike you I did not link this training to any particular culture like Tibetan or Native American, although somebody I was involved with at that time was adept in the native American tradition. But I did not identify with anything in particular, like you I was very much `myself'. Most initiated adepts, so far as I can tell, are drawn at some stage to Tibet and the misty hills of India. This does indeed seem to be a primary training ground, and it is easy to see why given the emphasis on self discipline, meditation techniques etc, it would be pretty hard to deny the efficacy of method and also to deny that the Buddhist stream (along with the others) ran into the Christ stream at the turning point of history. You should be aware, though, that just as Kabbalah has its Qlippoth, Sufism has its Islamic brotherhoods and Christianity has its inquisitors, Tibetan Buddhism has its demonology. When you start thinking in terms of minor gods rather than God you are open to trouble as you probably don't know what you are dealing with but are focusing instead on amassing power. Remember that during and prior to WWII it was to Tibet that the Nazis turned for the dark power that did so much to twist and subvert the magial aether, to pollute the magical streams. Remember as well that there are two eternities, one attained via dissolution and resolution in Christ, the other via crystalisation of an egoic diamond body. As an aside, I agree that so much of the `demonic manifestation' associated with the lesser guardian is actually ego, outward expressions of our vices and weaknesses etc -- but none the less terrifying for it. I had a 3.5 month trial with it at the end of last year and it all but did me in, I even asked for help, although the act of asking and gaining moral support -- acknowledging the reality of the problem -- was enough to help see it off. I was astrally attacked every single night, in the end I had no energy even to be afraid. The sign of the cross banished it.

\hfill

\texttt{Caleb Cooper on 2012-01-03 at 15:45 said:}

Actually, the training I did was core-shamanism, which is not based in any particular tradition. It compares all shamanic traditions and extracts its principles from their common elements. I use different practices that I consider in harmony with Tradition, though the aegis I operate under in Hermetic Christianity, which I consider superior, or rather, Heaven was more fully fulfilled on earth than it ever has been in the person of Christ. Which is why I've started the process of entering into communion with Rome. When I went to my first mass I was surprised just how much power was present, even in its attenuated form. Very healing for my soul.

``When you start thinking in terms of minor gods rather than God you are open to trouble as you probably don't know what you are dealing with but are focusing instead on amassing power.''

I generally agree here. When you come into resonance with the God behind the gods you don't typically have much need for the latter. When I seek a god, it's not rather than God, but as an aspect of God, like a facet of one jewel. Sometimes God is too big and blinding for a particular work, like trying to use a diamond sledgehammer when a jeweler's chisel would be more appropriate.

Specifically, for the work I was doing I needed an image at the same frequencies as the solar chakra. God contains all frequencies. You could extract the specific frequency from God, but at that point you have something that's less than God, so if it appears to you it will be as an aspect; the diamon of a regal lion if approached shamanically, perhaps as the Hindu god Rudra, or in my case when I went back and successfully completed the trial the Sun as symbol. So I guess this proves your point. If I'd stayed within the aegis of Hermetic astrology rather than stumbling into the realm of gods the first time I probably wouldn't have ran into so much trouble:)

\hfill

\texttt{Charlotte Cowell on 2012-01-03 at 16:07 said:}

funnily enough I had my `first' communion yesterday, well second -- I did try it earlier this year but found it genuinely appalling and very frightening as the church was wrong... .this time I sought out an authentic Latin mass and it lasted 2 hours but I found very soothing, it was a relief not to have to do any `work', but just let it all happen. I also felt it had a profound healing effect, although it's likely I'll always have an inclination to be independent in some way, my communions have always tended to in spirit or within the secret church and school.

I was actually converted to Christianity by a Benedictine 16 years ago but it's taken me this long to come around. I felt Catholic in terms of orientation and belief, but feared the communion due to issues with the Churh organisation and hierarchy -- fear of priests, basically.

\hfill

\texttt{Charlotte Cowell on 2012-01-03 at 16:08 said:}

PS for solar influences, why not Apollo? This also leads to profound insights about Orion and Artemis. Not to mention Dionysus etc... .


\hfill

\end{sffamily}
\end{footnotesize}

